\documentclass[a4paper]{arrowhead}

\usepackage[yyyymmdd]{datetime}
\usepackage{etoolbox}
\usepackage[utf8]{inputenc}
\usepackage[T1]{fontenc}
\usepackage{multirow}
\usepackage{longtable}
\usepackage{tikz}
\usetikzlibrary{positioning,shapes.multipart,arrows.meta}

\renewcommand{\dateseparator}{-}
\setlength{\emergencystretch}{3em}
\lstset{language=HTTP,breakatwhitespace=false,columns=fullflexible}

\newcommand{\fparam}[1]{\textit{\textcolor{ArrowheadBlue}{#1}}}

%% Special references
\newcommand{\mref}[2][]{{\textcolor{ArrowheadPurple}{\hyperref[sec:model:#2]{\ifstrempty{#1}{#2}{#1}}}}}
\newcommand{\oref}[1]{{\textcolor{ArrowheadBlue}{\hyperref[sec:operations:#1]{#1}}}}
\newcommand{\pdef}[1]{{\textcolor{ArrowheadGrey}{#1 \label{sec:model:primitives:#1} \label{sec:model:primitives:#1s}}}}
\newcommand{\pref}[1]{{\textcolor{ArrowheadGrey}{\hyperref[sec:model:primitives:#1]{#1}}}}

% as template_idd.tex; the input and output arguments are passed already formatted (\mref or \pref)
\newrobustcmd\osubsection[5]{
  \addtocounter{subsection}{1}
  \addcontentsline{toc}{subsection}{\protect\numberline{\thesubsection}operation \textcolor{ArrowheadBlue}{#1}}
  \subsection*{
    \thesubsection\quad
    #2 \textcolor{ArrowheadPurple}{#3} \\
    \small
    \hspace*{0.075\textwidth}\begin{minipage}{0.1\textwidth}
      \vspace*{1mm}
      Operation: \\
      \notblank{#4}{Input: \\}{}
      \notblank{#5}{Output: \\}{}
    \end{minipage}
    \begin{minipage}{0.825\textwidth}
      \vspace*{1mm}
      \textcolor{ArrowheadBlue}{#1} \\
      \notblank{#4}{#4 \\}{}
      \notblank{#5}{#5 \\}{}
    \end{minipage}
  }
  \label{sec:operations:#1}
}
\newrobustcmd\msubsection[2]{
  \addtocounter{subsection}{1}
  \addcontentsline{toc}{subsection}{\protect\numberline{\thesubsection}#1 \textcolor{ArrowheadPurple}{#2}}
  \subsection*{\thesubsection\quad#1 \textcolor{ArrowheadPurple}{#2}}
  \label{sec:model:#2} \label{sec:model:#2s}
}
%%

\begin{document}
\sloppy

%% Arrowhead Document Properties
\ArrowheadTitle{onboardingManagement HTTP/TLS/JSON}
\ArrowheadServiceID{onboardingManagement}
\ArrowheadType{Interface Design Description}
\ArrowheadTypeShort{IDD}
\ArrowheadVersion{5.2.2}
\ArrowheadDate{2026-10-30}
\ArrowheadAuthor{Jerker Delsing}
\ArrowheadStatus{RELEASE CANDIDATE}
\ArrowheadContact{jerker.delsing@ltu.se}
\ArrowheadFooter{\href{www.arrowhead.eu}{www.arrowhead.eu}}
\ArrowheadSetup
%%

%% Front Page
\begin{center}
  \vspace*{1cm}
  \huge{\arrowtitle}

  \vspace*{0.2cm}
  \LARGE{\arrowtype}
  \vspace*{1cm}

  \Large{Service ID: \textit{"\arrowid"}}

  \vspace*{\fill}

  % Front Page Image

  \vspace*{1cm}
  \vspace*{\fill}

  % Front Page Abstract
  \begin{abstract}
    This document describes the HTTP/1.1 protocol with TLS payload security and JSON payload encoding variant of the \texttt{onboarding\allowbreak{}Management} service of the Onboarding system.
  \end{abstract}

  \vspace*{1cm}
\end{center}

\newpage
%%

%% Table of Contents
\tableofcontents
\newpage
%%

\section{Overview}
\label{sec:overview}

This document describes the \texttt{onboarding\allowbreak{}Management} service interface of the Onboarding system, which provides operator approval and management of device enrollments.

It's implemented using protocol, encoding as stated in below table:
\begin{table}[ht!]
  \centering
  \begin{tabular}{|l|l|l|}
    \hline
    \rowcolor{gray!33} Profile type & Type & Version \\ \hline
    Transfer protocol & HTTP & 1.1 \cite{rfc9112} \\ \hline
    Data encryption & TLS & 1.3 \cite{rfc8446}, 1.2 \cite{rfc5246} \\ \hline
    Encoding & JSON & RFC 8259 \cite{rfc8259} \\ \hline
    Compression & N/A & - \\ \hline
    Semantics & N/A & - \\ \hline
    Ontology & N/A & - \\ \hline
  \end{tabular}
  \caption{Communication and semantics details used for the \texttt{onboarding\allowbreak{}Management} service interface}
  \label{tab:comunication_semantics_profile}
\end{table}

This document provides the Interface Design Description IDD to the \textit{onboardingManagement -- Service Description} document \cite{onb-onboardingManagement-sd}.
For further details about how this service is meant to be used, please consult that document.

The rest of this document describes how to realize the \texttt{onboarding\allowbreak{}Management} service
interface in detail. Both in terms of its operations (Section \ref{sec:operations}) and its information model (Section \ref{sec:model}).

\subsection{Interface registration}

The service is registered in the ServiceRegistry with the interface template \texttt{generic\_\allowbreak{}https}, policy \texttt{CERT\_\allowbreak{}AUTH}, access address \texttt{onboarding}, port 8435 and base path \texttt{/\allowbreak{}onboarding/\allowbreak{}mgmt}. The \texttt{operations} property of the interface lists the operations below with their method and path relative to the base path (path parameters are placeholders such as \texttt{\{id\}} that the consumer fills in):

\noindent\begin{longtable}{| p{5.2cm} | p{1.6cm} | p{8.4cm} |} \hline
\rowcolor{gray!33} Operation & Method & Path \\ \hline \endhead
\texttt{query-\allowbreak{}enrollments} & GET & \texttt{/\allowbreak{}enrollments} \\ \hline
\texttt{get-\allowbreak{}enrollment} & GET & \texttt{/\allowbreak{}enrollments/\allowbreak{}\{device\allowbreak{}Name\}} \\ \hline
\texttt{approve-\allowbreak{}device} & PUT & \texttt{/\allowbreak{}enrollments/\allowbreak{}\{device\allowbreak{}Name\}} \\ \hline
\texttt{offboard-\allowbreak{}device} & POST & \texttt{/\allowbreak{}enrollments/\allowbreak{}\{device\allowbreak{}Name\}/\allowbreak{}offboard} \\ \hline
\texttt{remove-\allowbreak{}enrollment} & DELETE & \texttt{/\allowbreak{}enrollments/\allowbreak{}\{device\allowbreak{}Name\}} \\ \hline
\end{longtable}

\newpage

\section{Service Operations}
\label{sec:operations}

The interfaces of the \texttt{onboarding\allowbreak{}Management} service, its operations, data models and implementation are provided below.
In particular, each subsection first names the HTTP method and path used to call the operation, after which it names the abstract operation from the \texttt{onboarding\allowbreak{}Management} SD document, as well as input and output types.

All operations require the requester's identity, presented according to the provider's authentication policy: the client certificate (mutual TLS) and, with the \texttt{outsourced} policy, \texttt{Authorization: Bearer IDENTITY-TOKEN//<token>}; systems holding the Authentication secret key send \texttt{Bearer AUTHENTICATOR-KEY//<system>//<hash>}, and with the \texttt{declared} policy \texttt{Bearer SYSTEM//<system>} is sent. Request and response bodies are \texttt{application/json} unless stated otherwise.
Errors are answered with an \mref{ErrorMessage} body and the status codes \texttt{400 Bad Request} (invalid parameter), \texttt{401 Unauthorized} (missing or invalid identity), \texttt{403 Forbidden} (not permitted, e.g. by the management, authorization or blacklist filter), \texttt{404 Not Found}, \texttt{408 Request Timeout}, \texttt{423 Locked}, \texttt{500 Internal Server Error} and \texttt{503 Service Unavailable} (a system it depends on failed).

\osubsection{query-enrollments}{GET}{/\allowbreak{}onboarding/\allowbreak{}mgmt/\allowbreak{}enrollments}{}{\mref{EnrollmentList}}

Returns all enrollments.

A successful call returns \texttt{200 OK} with the output in the body.

\begin{lstlisting}[label={lst:query-enrollments},caption={An \oref{query-enrollments} invocation.}]
GET /onboarding/mgmt/enrollments HTTP/1.1
Host: onboarding:8435
Authorization: Bearer IDENTITY-TOKEN//<token>
\end{lstlisting}

\begin{lstlisting}[label={lst:query-enrollments-resp},caption={An \oref{query-enrollments} response.}]
{
  "entries": [
    {
      "deviceName": "Raspberry1",
      "bootstrapFingerprint": "9f86d081884c7d659a2feaa0c55ad015a3bf4f1b2b0b822cd15d6c15b0f00a08",
      "approvedSystems": [
        "string"
      ],
      "state": "USED",
      "deviceCertificateId": 1,
      "deviceCertificateExpiresAt": "2026-10-22T12:00:00Z",
      "systems": [
        {
          "systemName": "TemperatureSensor",
          "certificateId": 1,
          "certificateExpiresAt": "2026-10-22T12:00:00Z",
          "issuedAt": "2026-10-22T12:00:00Z"
        }
      ],
      "approvedBy": "string",
      "approvedAt": "2026-10-22T12:00:00Z",
      "updatedAt": "2026-10-22T12:00:00Z",
      "lastMessage": "string"
    }
  ],
  "count": 1
}
\end{lstlisting}

\osubsection{get-enrollment}{GET}{/\allowbreak{}onboarding/\allowbreak{}mgmt/\allowbreak{}enrollments/\allowbreak{}\fparam{\{deviceName\}}}{\pref{String}}{\mref{EnrollmentView}}

Returns the enrollment of a device.

Path parameter: \fparam{deviceName}.

A successful call returns \texttt{200 OK} with the output in the body.

\begin{lstlisting}[label={lst:get-enrollment},caption={An \oref{get-enrollment} invocation.}]
GET /onboarding/mgmt/enrollments/<deviceName> HTTP/1.1
Host: onboarding:8435
Authorization: Bearer IDENTITY-TOKEN//<token>
\end{lstlisting}

\begin{lstlisting}[label={lst:get-enrollment-resp},caption={An \oref{get-enrollment} response.}]
{
  "deviceName": "Raspberry1",
  "bootstrapFingerprint": "9f86d081884c7d659a2feaa0c55ad015a3bf4f1b2b0b822cd15d6c15b0f00a08",
  "approvedSystems": [
    "string"
  ],
  "state": "USED",
  "deviceCertificateId": 1,
  "deviceCertificateExpiresAt": "2026-10-22T12:00:00Z",
  "systems": [
    {
      "systemName": "TemperatureSensor",
      "certificateId": 1,
      "certificateExpiresAt": "2026-10-22T12:00:00Z",
      "issuedAt": "2026-10-22T12:00:00Z"
    }
  ],
  "approvedBy": "string",
  "approvedAt": "2026-10-22T12:00:00Z",
  "updatedAt": "2026-10-22T12:00:00Z",
  "lastMessage": "string"
}
\end{lstlisting}

\osubsection{approve-device}{PUT}{/\allowbreak{}onboarding/\allowbreak{}mgmt/\allowbreak{}enrollments/\allowbreak{}\fparam{\{deviceName\}}}{\mref{EnrollmentRequest}}{\mref{EnrollmentView}}

Approves a device: bootstrap-certificate fingerprint and the systems it may onboard.

Path parameter: \fparam{deviceName}.

A successful call returns \texttt{200 OK} with the output in the body.

\begin{lstlisting}[label={lst:approve-device},caption={An \oref{approve-device} invocation.}]
PUT /onboarding/mgmt/enrollments/<deviceName> HTTP/1.1
Host: onboarding:8435
Authorization: Bearer IDENTITY-TOKEN//<token>
Content-Type: application/json

{
  "bootstrapFingerprint": "9f86d081884c7d659a2feaa0c55ad015a3bf4f1b2b0b822cd15d6c15b0f00a08",
  "systems": [
    "string"
  ]
}
\end{lstlisting}

\begin{lstlisting}[label={lst:approve-device-resp},caption={An \oref{approve-device} response.}]
{
  "deviceName": "Raspberry1",
  "bootstrapFingerprint": "9f86d081884c7d659a2feaa0c55ad015a3bf4f1b2b0b822cd15d6c15b0f00a08",
  "approvedSystems": [
    "string"
  ],
  "state": "USED",
  "deviceCertificateId": 1,
  "deviceCertificateExpiresAt": "2026-10-22T12:00:00Z",
  "systems": [
    {
      "systemName": "TemperatureSensor",
      "certificateId": 1,
      "certificateExpiresAt": "2026-10-22T12:00:00Z",
      "issuedAt": "2026-10-22T12:00:00Z"
    }
  ],
  "approvedBy": "string",
  "approvedAt": "2026-10-22T12:00:00Z",
  "updatedAt": "2026-10-22T12:00:00Z",
  "lastMessage": "string"
}
\end{lstlisting}

\osubsection{offboard-device}{POST}{/\allowbreak{}onboarding/\allowbreak{}mgmt/\allowbreak{}enrollments/\allowbreak{}\fparam{\{deviceName\}}/\allowbreak{}offboard}{\pref{String}}{\mref{OffboardingReport}}

Off-boards a device on behalf of the operator (as \texttt{onboarding.\allowbreak{}offboard-\allowbreak{}device}).

Path parameter: \fparam{deviceName}.

A successful call returns \texttt{200 OK} with the output in the body.

\begin{lstlisting}[label={lst:offboard-device},caption={An \oref{offboard-device} invocation.}]
POST /onboarding/mgmt/enrollments/<deviceName>/offboard HTTP/1.1
Host: onboarding:8435
Authorization: Bearer IDENTITY-TOKEN//<token>
\end{lstlisting}

\begin{lstlisting}[label={lst:offboard-device-resp},caption={An \oref{offboard-device} response.}]
{
  "deviceName": "Raspberry1",
  "completed": [
    "string"
  ],
  "failed": [
    "string"
  ]
}
\end{lstlisting}

\osubsection{remove-enrollment}{DELETE}{/\allowbreak{}onboarding/\allowbreak{}mgmt/\allowbreak{}enrollments/\allowbreak{}\fparam{\{deviceName\}}}{\pref{String}}{\mref{OffboardingReport}}

Off-boards the device and removes its enrollment.

Path parameter: \fparam{deviceName}.

A successful call returns \texttt{200 OK} with the output in the body.

\begin{lstlisting}[label={lst:remove-enrollment},caption={An \oref{remove-enrollment} invocation.}]
DELETE /onboarding/mgmt/enrollments/<deviceName> HTTP/1.1
Host: onboarding:8435
Authorization: Bearer IDENTITY-TOKEN//<token>
\end{lstlisting}

\begin{lstlisting}[label={lst:remove-enrollment-resp},caption={An \oref{remove-enrollment} response.}]
{
  "deviceName": "Raspberry1",
  "completed": [
    "string"
  ],
  "failed": [
    "string"
  ]
}
\end{lstlisting}

\newpage

\section{Data Models}
\label{sec:model}

Here, all data objects that can be part of the service calls associated with this service are listed in alphabetic order.
Note that each subsection, which describes one type of object, begins with the \textit{struct} keyword, which is meant to denote a JSON \pref{Object} that must contain certain fields, or names, with values conforming to explicitly named types.
Fields that are not needed may be omitted or \texttt{null} in requests; empty fields are omitted in responses.
As a complement to the primary types defined in this section, there is also a list of secondary types in Section \ref{sec:model:primitives}, which are used to represent things like hashes, identifiers and texts.

\msubsection{struct}{EnrollmentList}

List of enrollments.


\noindent\begin{tabularx}{\textwidth}{| >{\raggedright\arraybackslash}p{3.9cm} | >{\raggedright\arraybackslash\small}p{5.3cm} | X |} \hline
\rowcolor{gray!33} Object Field & Value Type & Description \\ \hline
"entries" & \pref{Array}$<$\allowbreak{}\mref{EnrollmentView}$>$ & Matching records of the current page. \\ \hline
"count" & \pref{Number} & Number of enrollments. \\ \hline
\end{tabularx}



\msubsection{struct}{EnrollmentRequest}

Operator approval of a device.


\noindent\begin{tabularx}{\textwidth}{| >{\raggedright\arraybackslash}p{3.9cm} | >{\raggedright\arraybackslash\small}p{5.3cm} | X |} \hline
\rowcolor{gray!33} Object Field & Value Type & Description \\ \hline
"bootstrap\allowbreak{}Fingerprint" & \pref{String} & SHA-256 fingerprint of the approved bootstrap certificate. \\ \hline
"systems" & \pref{Array}$<$\allowbreak{}\pref{String}$>$ & Systems the device may onboard. \\ \hline
\end{tabularx}



\msubsection{struct}{EnrollmentView}

An enrollment (approved device).


\noindent\begin{tabularx}{\textwidth}{| >{\raggedright\arraybackslash}p{3.9cm} | >{\raggedright\arraybackslash\small}p{5.3cm} | X |} \hline
\rowcolor{gray!33} Object Field & Value Type & Description \\ \hline
"device\allowbreak{}Name" & \pref{String} & Name of the device hosting the system. \\ \hline
"bootstrap\allowbreak{}Fingerprint" & \pref{String} & SHA-256 fingerprint of the approved bootstrap certificate. \\ \hline
"approved\allowbreak{}Systems" & \pref{Array}$<$\allowbreak{}\pref{String}$>$ & Systems the operator approved for the device. \\ \hline
"state" & \pref{String} & Enrollment state: \texttt{APPROVED}, \texttt{DEVICE\_\allowbreak{}REGISTERED} or \texttt{OFFBOARDED}. \\ \hline
"device\allowbreak{}Certificate\allowbreak{}Id" & \pref{Number} & Certificate\allowbreak{}Authority record of the device certificate. \\ \hline
"device\allowbreak{}Certificate\allowbreak{}Expires\allowbreak{}At" & \pref{DateTime} & Expiry of the device certificate. \\ \hline
"systems" & \pref{Array}$<$\allowbreak{}\mref{SystemEnrollmentView}$>$ & Systems onboarded for the device. \\ \hline
"approved\allowbreak{}By" & \pref{String} & Operator who approved the device. \\ \hline
"approved\allowbreak{}At" & \pref{DateTime} & Time of the approval. \\ \hline
"updated\allowbreak{}At" & \pref{DateTime} & UTC time of the last change. \\ \hline
"last\allowbreak{}Message" & \pref{String} & Last status message of the enrollment. \\ \hline
\end{tabularx}



\msubsection{struct}{ErrorMessage}




\noindent\begin{tabularx}{\textwidth}{| >{\raggedright\arraybackslash}p{3.9cm} | >{\raggedright\arraybackslash\small}p{5.3cm} | X |} \hline
\rowcolor{gray!33} Object Field & Value Type & Description \\ \hline
"error\allowbreak{}Message" & \pref{String} & Error text, if the event is an error. \\ \hline
"error\allowbreak{}Code" & \pref{Number} &  \\ \hline
"exception\allowbreak{}Type" & \pref{ExceptionType} &  \\ \hline
"origin" & \pref{String} &  \\ \hline
\end{tabularx}



\msubsection{struct}{OffboardingReport}

Outcome of an off-boarding.


\noindent\begin{tabularx}{\textwidth}{| >{\raggedright\arraybackslash}p{3.9cm} | >{\raggedright\arraybackslash\small}p{5.3cm} | X |} \hline
\rowcolor{gray!33} Object Field & Value Type & Description \\ \hline
"device\allowbreak{}Name" & \pref{String} & The off-boarded device. \\ \hline
"completed" & \pref{Array}$<$\allowbreak{}\pref{String}$>$ & Off-boarding steps that succeeded. \\ \hline
"failed" & \pref{Array}$<$\allowbreak{}\pref{String}$>$ & Off-boarding steps that failed. \\ \hline
\end{tabularx}



\msubsection{struct}{SystemEnrollmentView}

A system onboarded for a device.


\noindent\begin{tabularx}{\textwidth}{| >{\raggedright\arraybackslash}p{3.9cm} | >{\raggedright\arraybackslash\small}p{5.3cm} | X |} \hline
\rowcolor{gray!33} Object Field & Value Type & Description \\ \hline
"system\allowbreak{}Name" & \pref{String} & Name of the system, PascalCase (e.g. \texttt{TemperatureSensor}). \\ \hline
"certificate\allowbreak{}Id" & \pref{Number} & Certificate\allowbreak{}Authority record of the system certificate. \\ \hline
"certificate\allowbreak{}Expires\allowbreak{}At" & \pref{DateTime} & Expiry of the system certificate. \\ \hline
"issued\allowbreak{}At" & \pref{DateTime} & Time the certificate was issued. \\ \hline
\end{tabularx}



\subsection{Primitives}
\label{sec:model:primitives}

As all messages are encoded using the JSON format \cite{rfc8259}, the following primitive constructs, part of that standard, become available.
Note that the official standard is defined in terms of parsing rules, while this list only concerns syntactic information.
Furthermore, the \pref{Object} and \pref{Array} types are given optional generic type parameters, which are used in this document to signify when pair values or elements are expected to conform to certain types.

\noindent\begin{tabularx}{\textwidth}{| p{4.6cm} | X |} \hline
\rowcolor{gray!33} Type & Description \\ \hline
\pdef{Array} & A JSON Array \cite{rfc8259}; \pref{Array}$<$A$>$ has elements of type A. \\ \hline
\pdef{DateTime} & JSON String with an ISO 8601 UTC timestamp, e.g. \texttt{2026-10-22T12:00:00Z}. \\ \hline
\pdef{ExceptionType} & One of \texttt{ARROWHEAD}, \texttt{INVALID\_\allowbreak{}PARAMETER}, \texttt{AUTH}, \texttt{FORBIDDEN}, \texttt{DATA\_\allowbreak{}NOT\_\allowbreak{}FOUND}, \texttt{TIMEOUT}, \texttt{LOCKED}, \texttt{INTERNAL\_\allowbreak{}SERVER\_\allowbreak{}ERROR}, \texttt{EXTERNAL\_\allowbreak{}SERVER\_\allowbreak{}ERROR} (JSON String). \\ \hline
\pdef{Number} & JSON Number, an integer. \\ \hline
\pdef{Object} & JSON Object; \pref{Object}$<$A$>$ has string keys and values of type A. \\ \hline
\pdef{String} & JSON String. \\ \hline
\end{tabularx}

\newpage

\bibliographystyle{IEEEtran}
\bibliography{bibliography}

\newpage

\section{Revision History}
\subsection{Amendments}

\noindent\begin{tabularx}{\textwidth}{| p{1cm} | p{3cm} | p{2cm} | X | p{4cm} |} \hline
\rowcolor{gray!33} No. & Date & Version & Subject of Amendments & Author \\ \hline

1 & 2026-10-22 & \arrowversion & Document created from the Arrowhead LaTeX templates for the 5.2.2 Java Spring implementation, including the local extensions of this workspace. & Jerker Delsing \\ \hline
2 & 2026-10-30 & \arrowversion & Checked against the primary Java source tree (Eclipse-Arrowhead-java): services, operations, paths and data models unchanged; missing primitive definitions added. & Claude Opus 5.5 (draft) \\ \hline
\end{tabularx}

\subsection{Quality Assurance}

\noindent\begin{tabularx}{\textwidth}{| p{1cm} | p{3cm} | p{2cm} | X |} \hline
\rowcolor{gray!33} No. & Date & Version & Approved by \\ \hline

1 & & \arrowversion  &  \\ \hline

\end{tabularx}

\end{document}
